\chapter{Evaluation and Discussion}
\label{ch:evaluation}

\section{Central Empirical Finding}

The strongest result from the current corpus is not a detection F1 score. It is the systematic mismatch between benign folder labels and rule-based triage signals.

In \texttt{dataset\_manifest.csv} ($n{=}30$), \textbf{14 of 15} benign-labelled (\texttt{NoVirus}) rows are marked \texttt{uncertain=True}. Reasons include HIGH or CRITICAL verdicts, elevated process scores, injection counts, and C2-style connection counts despite benign labels. Only the W32.MyDoom.A.-NoVirus row is not uncertain. Malware-labelled rows are not flagged uncertain in the manifest.

Figure~\ref{fig:uncertainty_benign} summarises the split. Figure~\ref{fig:dataset_overview} places it in the wider label balance. This pattern means folder naming is an imperfect proxy for ground truth: ``clean'' baselines in memory forensics can still contain suspicious artefacts from security tools, prior activity, or capture timing.

\begin{figure}[htbp]
\centering
\includegraphics[width=0.62\textwidth]{fig_uncertainty_benign.pdf}
\caption{Benign-labelled samples: uncertain vs not uncertain ($n{=}15$).}
\label{fig:uncertainty_benign}
\end{figure}

For machine learning, uncertain benign rows should be excluded, down-weighted, or reported separately in evaluation. Treating all NoVirus folders as definitively benign will inflate false conclusions about separability.

\section{Proof-of-Concept Scope}

This chapter evaluates whether the pipeline produces consistent artefacts, meaningful graphs, and auditable evidence on a small research corpus. It does \textbf{not} claim production malware detection performance. Quantitative ML tables use placeholders until grouped cross-validation is run and checkpoints are archived (see \S\ref{sec:quant_ml_eval}).

\section{Descriptive Corpus Statistics}

Table~\ref{tab:corpus_summary} summarises manifest-level scale statistics (from \texttt{figures/dataset\_stats.json}). Graph size spans more than two orders of magnitude; both labels include very large graphs.

\begin{table}[htbp]
\centering
\caption{Corpus summary statistics (30 manifest rows).}
\label{tab:corpus_summary}
\begin{tabular}{lrrrr}
\toprule
\textbf{Statistic} & \textbf{Min} & \textbf{Median} & \textbf{Max} & \textbf{Notes} \\
\midrule
Nodes (all) & 86 & 18,036 & 27,362 & MyDoom pair smallest \\
Edges (all) & 7 & 21,544 & 32,556 & \\
Malware-labelled & 15 & --- & --- & No uncertain flags \\
Benign-labelled & 15 & --- & --- & 14 uncertain \\
\bottomrule
\end{tabular}
\end{table}

\begin{figure}[htbp]
\centering
\includegraphics[width=0.95\textwidth]{fig_dataset_overview.pdf}
\caption{Dataset composition: malware vs benign labels and uncertain benign subset.}
\label{fig:dataset_overview}
\end{figure}

\begin{figure}[htbp]
\centering
\includegraphics[width=0.95\textwidth]{fig_graph_scale.pdf}
\caption{Node and edge counts by label (log scale).}
\label{fig:graph_scale}
\end{figure}

All listed rows completed filter, graph, and analysis stages (\texttt{filter\_ok}, \texttt{graph\_ok}, \texttt{analyze\_ok}), giving a stable artefact base for discussion.

\section{Rule-Based Triage Analysis}

Rule-based signals in the manifest show false-positive pressure on NoVirus rows. Figure~\ref{fig:forensic_signals} compares mean injection count, C2 connections, and maximum process score by label. Malware-labelled means are not uniformly higher than benign-labelled means for every signal.

\begin{figure}[htbp]
\centering
\includegraphics[width=0.85\textwidth]{fig_forensic_signals.pdf}
\caption{Mean forensic signal values by label (manifest-derived).}
\label{fig:forensic_signals}
\end{figure}

Figure~\ref{fig:verdict_dist} shows that most samples receive HIGH or CRITICAL verdicts. Only the MyDoom pair reports LOW verdicts. Benign-labelled rows therefore often receive the same verdict bands as malware-labelled rows, which explains the uncertain flags.

\begin{figure}[htbp]
\centering
\includegraphics[width=0.78\textwidth]{fig_verdict_distribution.pdf}
\caption{Rule-based triage verdict categories across the corpus.}
\label{fig:verdict_dist}
\end{figure}

Figure~\ref{fig:family_coverage} shows each malware family contributes paired WithVirus and NoVirus rows. Family-grouped splits are mandatory for honest evaluation \cite{sommer2010outside}.

\begin{figure}[htbp]
\centering
\includegraphics[width=0.72\textwidth]{fig_family_coverage.pdf}
\caption{Manifest rows per malware family.}
\label{fig:family_coverage}
\end{figure}

\section{Quantitative Machine Learning Evaluation}
\label{sec:quant_ml_eval}

Grouped evaluation is implemented in \texttt{train.py} (StratifiedGroupKFold by \texttt{family}), \texttt{train\_binary\_model.py}, and \texttt{scripts/tabular\_baselines.py}. Saved \texttt{results\_*.json} files were not present in the repository at writing time; Table~\ref{tab:gnn_comparison}--\ref{tab:cv_protocol} therefore show ``--'' placeholders. After running experiments, regenerate tables with:

\begin{verbatim}
python train.py extracted_data/dataset_manifest.csv --cv stratified_group \
  --group-by family --model gine --epochs 80 --seed 42
python scripts/tabular_baselines.py
python scripts/generate_eval_latex.py
\end{verbatim}

% Auto-generated by scripts/generate_eval_latex.py
% Tabular baselines only; Track 1 GNN CV not reported numerically.

\subsection{Tabular baselines (manifest graph-level features)}
\begin{table}[htbp]
\centering
\caption{Tabular baseline CV on \texttt{extracted\_data} manifest (\texttt{tabular\_baselines.py}). Reference only---primary ML results are Tracks 2/3.}
\label{tab:tabular_baselines}
\begin{tabular}{lcccc}
\toprule
\textbf{Model} & \textbf{Precision} & \textbf{Recall} & \textbf{F1} & \textbf{ROC-AUC} \\
\midrule
Logistic Regression & 0.500 & 0.467 & 0.469 & 0.489 \\
Random Forest & 0.687 & 0.667 & 0.631 & 0.711 \\
\bottomrule
\end{tabular}
\end{table}


\subsection{Evaluation figures (placeholders until predictions exist)}

When \texttt{outputs/predictions.csv} is produced from grouped CV, \texttt{scripts/generate\_eval\_latex.py} can render confusion-matrix and ROC PDFs (\texttt{fig\_confusion\_matrix.pdf}, \texttt{fig\_roc\_curve.pdf}). Until then, only manifest-derived figures (including Figure~\ref{fig:uncertainty_benign}) are included.

% Uncomment after running generate_eval_latex.py with predictions:
% \begin{figure}[htbp]
% \centering
% \includegraphics[width=0.55\textwidth]{fig_confusion_matrix.pdf}
% \caption{Confusion matrix for GINE (grouped CV).}
% \label{fig:confusion_matrix}
% \end{figure}
%
% \begin{figure}[htbp]
% \centering
% \includegraphics[width=0.55\textwidth]{fig_roc_curve.pdf}
% \caption{ROC curve for GINE (grouped CV).}
% \label{fig:roc_curve}
% \end{figure}

A natural comparison is tabular classifiers on \texttt{graph\_attr} features versus full-graph GNNs on the same manifest version. If tabular models match GNNs, relational structure may add little on this small, heuristic-rich corpus; if GNNs improve under grouped CV, that supports the graph design.

\section{Failure Cases and Edge Families}

\subsection{W32.MyDoom.A. pair (very small graphs)}

W32.MyDoom.A.-NoVirus (86 nodes) and W32.MyDoom.A.-WithVirus (164 nodes) are outliers. Both report LOW verdicts and zero injections; neither benign row is uncertain. Models trained mainly on large graphs may underfit or overfit on tiny topologies---pooling and normalisation should be checked on this pair.

\subsection{Cerber pair (similar scale, divergent C2 counts)}

Cerber-WithVirus (7,904 nodes) and Cerber-NoVirus (9,382 nodes) have comparable scale. The malware row reports 24 C2 connections and \texttt{lolbin\_c2\_found=True}; the benign row reports 7 C2 connections and is uncertain. Both share HIGH verdicts and six injections in the manifest. A GNN may not separate these folders if features emphasise shared injection structure more than family-specific infection cues.

\subsection{GoldenEye pair (high injection on both rows)}

GoldenEye-WithVirus has the highest injection count (41). GoldenEye-NoVirus still reports 8 injections, 13 hidden processes, and a CRITICAL verdict. This is among the clearest uncertain-benign examples: triage treats both snapshots as highly suspicious.

\subsection{WannaCry case study}

WannaCry-WithVirus (8,098 nodes) shows injection-related attack-chain steps and ransom-note signals. WannaCry-NoVirus (7,590 nodes) reports nine injections, a HIGH verdict, and \texttt{uncertain=True}. Figure~\ref{fig:wannacry_nodes} shows memory regions dominate node types---typical for VAD-heavy graphs.

\begin{figure}[htbp]
\centering
\includegraphics[width=0.62\textwidth]{fig_wannacry_node_types.pdf}
\caption{Node-type composition for WannaCry-WithVirus.}
\label{fig:wannacry_nodes}
\end{figure}

\begin{table}[htbp]
\centering
\caption{WannaCry manifest comparison (selected fields).}
\label{tab:wannacry_compare}
\begin{tabular}{lrrrr}
\toprule
\textbf{Folder} & \textbf{Label} & \textbf{Nodes} & \textbf{Injections} & \textbf{Uncertain} \\
\midrule
WannaCry-WithVirus & Malware & 8,098 & 7 & No \\
WannaCry-NoVirus & Benign & 7,590 & 9 & Yes \\
\bottomrule
\end{tabular}
\end{table}

\subsection{When a GNN may not help}

GNNs add value when relational structure carries signal beyond graph-level heuristics. They may add little when: (1) labels are dominated by triage noise; (2) graphs differ mainly in size, not infection topology; (3) train and test families are too few for stable grouped metrics; or (4) edge and node features already encode the decision rules. The uncertain-benign prevalence pushes toward analyst-review workflows and two-model conformity scores rather than hard binary automation.

\section{Binary and Two-Model Paths}

The binary GNN path supports calibrated probabilities and ambiguous states \cite{guo2017calibration}. On noisy labels, a single threshold can encode triage noise rather than infection state. The two-model path compares malware and benign conformity \cite{ruff2018deep}, surfacing cases where both scores are elevated---consistent with uncertain benign rows.

Saved checkpoints were not archived in the repository at writing time. Implementation capability is documented; grouped metrics require the user to train, save \texttt{*.pt} files, and regenerate LaTeX tables.

\section{Explainability}

The project records top graph attributes, important nodes, and edge pairs in JSON evidence. This is not a complete explanation method \cite{doshi2017interpretable}, but it exceeds a single opaque score. GNNExplainer \cite{ying2019gnnexplainer} remains future work.

\section{What Results Prove and Do Not Prove}

\textbf{Proven on current artefacts:}
\begin{itemize}
\item Volatility artefacts can be converted into typed graphs from 86 to 27,362 nodes.
\item Rule-based triage produces HIGH/CRITICAL verdicts on most benign-labelled rows (14/15 uncertain).
\item Manifest governance supports reproducible, family-aware evaluation design.
\item Evidence JSON and uncertainty fields are implementable alongside GNN scores.
\end{itemize}

\textbf{Not proven:}
\begin{itemize}
\item Superior detection accuracy versus production AV or EDR on enterprise hosts.
\item Generalisation to unseen malware families without grouped CV on larger $n$.
\item That GNN structure beats strong tabular baselines on the same manifest (pending experiments).
\item Robustness to adversarial graph manipulation \cite{apruzzese2023sok}.
\end{itemize}

\section{Limitations}

\begin{itemize}
\item Small corpus ($n{=}30$) with high metric variance under grouped CV.
\item Noisy benign labels (14/15 uncertain).
\item Family-level pairing limits independent test families.
\item Capture timing and Windows build affect plugin output.
\item Heuristic graph attributes couple rules and learning.
\item Missing archived checkpoints and prediction CSVs at submission time.
\end{itemize}

Sommer and Paxson \cite{sommer2010outside} and Sculley et al.\ \cite{sculley2015debt} apply: evaluation design and pipeline debt can dominate headline metrics.

\section{Practical Value}

The pipeline supports evidence-led triage \cite{case2017memory}: organise plugin output, surface rule verdicts, attach GNN scores as a secondary signal, and export JSON for analyst review. The uncertain-benign finding is operationally relevant---teams should expect HIGH verdicts on baseline images without treating them automatically as labelling errors.

\section{Summary}

\begin{enumerate}
\item Benign-labelled rows frequently exhibit HIGH/CRITICAL triage signals (14/15 uncertain)---the central empirical finding.
\item Graph construction scales from tiny to very large memory snapshots.
\item Rule-based triage alone cannot define benign status from folder names.
\item Grouped CV and tabular baselines are implemented; numeric ML tables await experiment runs.
\item Structured evidence improves usefulness over a single probability.
\end{enumerate}
